\chapter{Dạy Con Bằng Nếp Sống}
\markboth{Dạy Con Bằng Nếp Sống}{Teach Children Through Daily Habits}

\section{Sao Bố Hay Tắt Điện?}
\begin{truyen}{Một Câu Hỏi Rất Nhỏ Nhưng Đủ Mở Ra Cả Một Bài Học}{A Small Question That Opens a Large Lesson}
\chuhoa{M}{ột buổi tối, đứa con hỏi: ``Sao bố hay tắt điện vậy?'' Người cha cười, rồi nhận ra đó không phải một câu hỏi về bóng đèn. Đó là câu hỏi về cách sống.}
\textit{(One evening, the child asked, ``Why do you always turn off the lights?'' The father smiled and realized that it was not really a question about a light bulb. It was a question about a way of living.)}

Anh bảo con rằng bố không sợ một bóng đèn, cũng không vì mấy nghìn đồng mà khắt khe. Điều bố sợ là sự cẩu thả lặp lại thành thói quen. Một người quen quên tắt điện, quên khóa vòi nước, quên cất đồ cho gọn, lâu ngày sẽ quen cả việc xem nhẹ những gì mình đang có.
\textit{(He told the child that he was not afraid of one light bulb, nor was he being strict over a tiny amount of money. What he feared was carelessness repeated until it became habit. A person who habitually forgets to turn off lights, close taps, or put things away may gradually become accustomed to taking everything he has for granted.)}

Nhiều bài học về tiền không bắt đầu từ đồng tiền. Chúng bắt đầu từ thái độ với những thứ rất nhỏ. Và trẻ con nhìn vào thái độ ấy còn kỹ hơn những lời dạy chính thức.
\textit{(Many lessons about money do not begin with money itself. They begin with one’s attitude toward small things. And children watch that attitude even more closely than formal lessons.)}
\end{truyen}

\begin{baihoc}
Muốn con quý tiền, trước hết phải dạy con quý những điều tiền mua được và những công sức đứng sau chúng.
\textit{(If you want children to value money, first teach them to value what money provides and the labor behind it.)}
\end{baihoc}

\begin{ghinhoanh}
Children learn values through details.

Small habits build big lessons.
\end{ghinhoanh}

\begin{tuhocnhanh}
1. Trẻ nhỏ đang nhìn thấy nếp sống nào ở em mỗi ngày? \textit{What daily habits are children seeing in you?}

2. Điều nhỏ nào em đang xem thường nhưng con em sẽ bắt chước rất nhanh? \textit{What small thing are you overlooking that your children may imitate quickly?}
\end{tuhocnhanh}
\ngancach

\section{Tăng Tử Giết Lợn}
\begin{truyen}{Lời Hứa Và Sự Nêu Gương}{Promise and Example}
\chuhoa{T}{ăng Tử từng giết lợn để giữ lời với con. Câu chuyện ấy không chỉ nói về chữ tín, mà còn cho thấy trẻ con học bằng điều mắt chúng thấy người lớn làm.}
\textit{(Zengzi once slaughtered a pig to keep a promise to his child. That story is not only about trustworthiness; it also shows that children learn from what they see adults do.)}

Trong việc sử dụng tiền cũng vậy. Cha mẹ nói rằng phải tiết kiệm nhưng lại tiêu theo cảm hứng, nói rằng phải biết đủ nhưng luôn chạy theo người khác, nói rằng phải quý lao động nhưng lại vứt bỏ thức ăn thừa, tất cả những mâu thuẫn ấy trẻ con đều nhìn ra.
\textit{(The same is true of money. Parents may say that one must save, but then spend according to mood; say that one must know enough, yet constantly chase others; say that labor should be valued, yet waste food. Children notice all these contradictions.)}

Người cha hiểu rằng muốn con lớn lên biết chắt chiu, anh phải bớt nói hay và sống thật hơn. Một bữa cơm ăn hết, một món đồ dùng bền, một lần từ chối mua thứ vô ích, đôi khi dạy mạnh hơn cả một bài giảng dài.
\textit{(The father understood that if he wanted his children to grow up frugal, he had to speak less beautifully and live more truthfully. Finishing a meal, using an item until it lasts, refusing a useless purchase: such acts often teach more strongly than a long lecture.)}
\end{truyen}

\begin{baihoc}
Trẻ con không học tiết kiệm từ khẩu hiệu. Chúng học từ dáng sống của cha mẹ.
\textit{(Children do not learn thrift from slogans. They learn it from the way their parents live.)}
\end{baihoc}

\begin{ghinhoanh}
Example teaches faster than words.

Children believe what adults live.
\end{ghinhoanh}

\begin{tuhocnhanh}
1. Điều gì em đang nói với con nhưng chính mình chưa làm được? \textit{What are you telling children that you yourself do not yet practice?}

2. Hôm nay em có thể nêu gương điều gì thật nhỏ nhưng thật rõ? \textit{What small but clear example can you set today?}
\end{tuhocnhanh}
\ngancach

\section{Cho Con Cùng Ghi Một Khoản Nhỏ}
\begin{truyen}{Dạy Con Phân Biệt Cần Và Muốn}{Teaching a Child to Distinguish Need and Want}
\chuhoa{C}{ó tối người cha gọi con lại, đặt cuốn sổ chi tiêu lên bàn và bảo con thử ghi cùng một vài khoản nhỏ trong ngày. Đứa trẻ ngạc nhiên vì chưa từng nghĩ người lớn cũng phải tính kỹ như vậy.}
\textit{(One evening, the father called his child over, placed the expense notebook on the table, and asked the child to help record a few small expenses from the day. The child was surprised, never having imagined that adults too had to calculate carefully.)}

Anh không làm việc ấy để con sợ tiền. Anh làm để con hiểu tiền có đường đi và đường về, và nếu không biết trông chừng thì nó mất rất nhanh. Anh hỏi con món nào là cần, món nào chỉ là thích, rồi để con tự nói ra. Đứa trẻ bắt đầu hiểu rằng không phải điều gì mình muốn cũng nên có ngay.
\textit{(He did not do this to make the child fear money. He did it so the child would understand that money has paths in and out, and if no one watches, it disappears quickly. He asked which items were needs and which were merely wants, letting the child answer. The child began to see that not everything desired should be owned at once.)}

Những buổi học như vậy rất giản dị. Nhưng về lâu dài, chúng giúp con lớn lên với cảm giác trân trọng hơn. Và đó là món quà đáng giá hơn nhiều món đồ chơi đắt tiền.
\textit{(Lessons like that are very simple. But in the long run, they help a child grow up with greater respect for things. And that is a gift worth more than many expensive toys.)}
\end{truyen}

\begin{baihoc}
Dạy con tiết kiệm không phải làm con thiếu thốn. Đó là giúp con có nội lực để sống bền hơn sau này.
\textit{(Teaching children thrift is not making them deprived. It is helping them develop the strength to live more steadily later in life.)}
\end{baihoc}

\begin{ghinhoanh}
Teach the child to choose.

Choice is the root of discipline.
\end{ghinhoanh}

\begin{tuhocnhanh}
1. Nếu dạy một đứa trẻ về tiền, em muốn dạy điều gì đầu tiên? \textit{If you were teaching a child about money, what would be the first lesson?}

2. Em có đang cho con quá nhiều để đổi lấy sự yên thân của người lớn không? \textit{Are you giving children too much just to buy adult convenience?}
\end{tuhocnhanh}

\begin{turanminh}
1. Tôi không được đòi con biết quý tiền nếu chính tôi còn sống dễ dãi với điện, nước, cơm áo và những khoản lặt vặt trong nhà.

2. Muốn con đỡ khổ về sau, tôi phải dạy con biết đủ từ những việc nhỏ hôm nay.

3. Mỗi lần tôi ăn hết bữa cơm, giữ đồ dùng bền và từ chối một món không cần là một lần tôi đang dạy con mà không cần lên giọng.
\end{turanminh}

\begin{dayhoc}
1. Chương này có thể dùng để nhắc học sinh rằng giáo dục không chỉ diễn ra trên lớp, mà đang diễn ra trong từng thói quen của gia đình.

2. Thầy có thể hỏi các em: ``Cha mẹ em đang dạy em điều gì qua cách dùng tiền mỗi ngày?''

3. Bài học nên nhấn là trẻ con học nhanh nhất từ điều mắt thấy, nên người lớn muốn dạy gì thì phải sống điều đó trước.
\end{dayhoc}